\documentclass[12pt]{article}

\usepackage{amsmath}
\usepackage[margin=1in]{geometry}
\usepackage{fancyhdr}
\usepackage{amsthm}
\newtheorem{lemma}{Lemma}
\usepackage{braket}
\usepackage{amssymb}
\usepackage{tikz}
\usetikzlibrary{shapes.geometric}
\usepackage{enumerate}
\usepackage[shortlabels]{enumitem}


\pagestyle{fancy}
\fancyhead[l]{Hudson Benites}
\fancyhead[c]{Homework \#7}
\fancyfoot[c]{\thepage}
\renewcommand{\headrulewidth}{0.2pt}
\setlength{\headheight}{15pt}
\fancyhead[r]{\today}

\begin{document}
\begin{enumerate}[start = 1, label = {\bfseries Question \arabic*:}, leftmargin=1in]
\item For the following proofs, suppose $F$ is a field, and $a,b\in F$. \begin{itemize}
    \item[(a)] If $a=0$ or $b=0$, we are done. If not, WLOG, let $a\neq 0$. Then, multiplying by $a^{-1}$ (which is guaranteed to exist by A13), we get $b=0$. Thus, $a=0$ or $b=0$.
    \item[(b)] Suppose $a^2=b^2$. Then, $a^2-b^2=(a-b)(a+b)=0$. Using the result from (a), $a-b=0$ or $a+b=0$. Thus, $a=b$ or $a=-b$.
    \item[] 
\end{itemize}
\item should be Q2 \begin{proof}
    Consider the set $L$, which is the set containing all lower bounds. It is bounded above by every element in $S$ (which exist because by assumption, $S$ is non-empty). So, $L$ has a supremum $l$.
    We claim that $l=\inf S$.\begin{description}
        \item[Lower Boundedness] By definition, $l$ is a lower bound.
        \item[Greatness] By definition, $l$ is the greater than every lower bound.
    \end{description} Thus, $l=\inf S$, and therefore $S$ admits an infimum.
\end{proof}
\item should be Q3 \begin{itemize}
    \item[(a)]\begin{proof}

\end{proof}
    \item[(b)] We claim that $b = \sup S$.
    \begin{description}
        \item[Upper Boundedness] Rewriting $S$ as $S=\Set{x\in \mathbb{R} | a \leq x < b}$, we can see that $b$ is an upper bound by definition.
        \item[Leastness] For the sake of contradiction, assume $b$ is not the supremum. Then, there exists a $\sup S = u$. By definition, we must have $a\leq u < b$. By Theorem 6.12 in the notes, we can construct
        a rational number $r$ with $u<r<b$. By extension, we have $a\leq u<r<b$. So, $r\in S$, and $u$ is not an upper bound. This a contradiction with our intial assumption, so $b$ must be the least upper bound.
    \end{description}
    Thus, $\sup S = b$.
\end{itemize}

\item should be q4. We compute that $\sup S = 1$ and $\inf S = 0$. We first prove the former.\begin{proof}
    The first element in $S$ is 1, and every succesive element is smaller, thus 1 is the max of $S$. So, $1=\sup S$.
\end{proof}
    Next, we prove $\inf S = 0$.\begin{proof}
        how tf u do ts i thought the rationals were incomplete
    \end{proof}

\item should be q5. \begin{proof}
    Suppose $\alpha$ is irrational, and $r\in \mathbb{Q}\setminus{0}$. Suppose for the sake of contradiction that $r\alpha \in \mathbb{Q}$. Then,
    we can write $r\alpha = \frac{c}{d}$ for $c,d\in \mathbb{Z}$. We can also write $r=\frac{a}{b}$ for $a,b\in \mathbb{Z}$. Then, we have $r\alpha = \frac{a}{b}\alpha= \frac{c}{d}$. Rearranging, we get that \begin{center}
        $\alpha = \frac{bc}{ad}$,
    \end{center} which means $\alpha \in \mathbb{Q}$ (since $a,b,c,d\in \mathbb{Z}$). However, this contradicts the assumption that $\alpha$ was irrational. Thus, $r\alpha$ is irrational for every $r\in \mathbb{Q}\setminus {0}$.
\end{proof}
\item should be q7
    \begin{itemize}
        \item[(a)] We compute $x_3$ in base 10 with this expression: \begin{center}
            $\sum\limits_{n=1}^{\infty}\frac{2}{3^{3n}}=\frac{2}{3}\sum\limits_{n=1}^{\infty}\frac{1}{3^n}=0.\overline{3}_{10}$.
        \end{center} It is in the Cantor set.
        \item[(b)] Similarly for $y_3$:\begin{center}
            $\sum\limits_{n=1}^{\infty}\frac{1}{3^{2n}}+\sum\limits_{n=1}^{\infty}\frac{1}{3^{3n}}=\frac{1}{9}\sum\limits_{n=1}^{\infty}\frac{1}{3^n}+\sum\limits_{n=1}^{\infty}\frac{1}{3^n}=\frac{4}{54}=0.0\overline{740}_{10}$.
        \end{center} you tell me if its in the cantor set twin
        \item[(c)] For $z_3$, there is only one digit to compute. It is in the third place to the right, so it is $\frac{1}{3^3}=\frac{1}{27}=0.0\overline{370}_{10}$
    \end{itemize}
\end{enumerate}
\end{document}